\documentclass[aspectratio=169,10pt]{beamer}
\usetheme[numbering=fraction,progressbar=frametitle]{metropolis}
\usepackage{fontspec}
\setsansfont{Latin Modern Sans}
\usepackage{amsmath,amssymb,mathtools,graphicx,booktabs}
\usepackage{polyglossia}
\setdefaultlanguage{english}
\setbeamertemplate{navigation symbols}{}
\setbeamerfont{frametitle}{size=\large}
\setbeamerfont{normal text}{size=\normalsize}
\setlength{\parskip}{.45em}
\hypersetup{pdfauthor={},pdfsubject={Lecture slides},pdfkeywords={},pdftitle={Sequences}}
\providecommand{\N}{\mathbb{N}}
\providecommand{\Z}{\mathbb{Z}}
\providecommand{\D}{\mathbb{D}}
\providecommand{\Q}{\mathbb{Q}}
\providecommand{\R}{\mathbb{R}}
\providecommand{\C}{\mathbb{C}}
\providecommand{\vect}[1]{\overrightarrow{#1}}
\providecommand{\vi}{\vec{\imath}}
\providecommand{\vj}{\vec{\jmath}}
\providecommand{\vk}{\vec{k}}
\providecommand{\norm}[1]{\left\lVert #1\right\rVert}
\providecommand{\vc}[2]{\begin{pmatrix}#1\\#2\end{pmatrix}}
\providecommand{\deter}[2]{\det\left(#1,#2\right)}
\providecommand{\abs}[1]{\left\lvert #1\right\rvert}
\providecommand{\dd}{\,\mathrm{d}}
\providecommand{\e}{\mathrm{e}}
\providecommand{\Card}{\operatorname{Card}}
\providecommand{\E}{\operatorname{E}}
\providecommand{\V}{\operatorname{V}}
\providecommand{\degre}{^{\circ}}
\providecommand{\un}[1]{\,\mathrm{#1}}
\providecommand{\sci}[2]{#1\times10^{#2}}
\providecommand{\pgcd}{\operatorname{PGCD}}
\providecommand{\ppcm}{\operatorname{PPCM}}
\title{Sequences}
\subtitle{Lecture and worked examples}
\date{}
\author{}
\begin{document}
\begin{frame}[plain]
\vfill
{\usebeamerfont{title}\inserttitle\par}
\medskip
{\usebeamerfont{subtitle}\insertsubtitle\par}
\medskip
{\usebeamercolor[fg]{progress bar}\rule{\textwidth}{.5pt}}
\vfill
\end{frame}
\begin{frame}{Core Course Notes}
A sequence \((u_n)\) is \textbf{arithmetic} when there is a real number \(r\) such that, for all \(n\), \(u_{n+1}=u_n+r\). The number \(r\) is the common difference. If the first term is \(u_0\), then \(u_n=u_0+nr\). If the first term is \(u_1\), then \(u_n=u_1+(n-1)r\).
\end{frame}
\begin{frame}{Core Course Notes}
To model an additive situation: 1) define \(u_n\) and state the unit; 2) identify the first term; 3) express the constant increase or decrease as the common difference \(r\); 4) write the general term; 5) answer with a sentence in context.
\end{frame}
\begin{frame}{Core Course Notes}
A shop sells \(480\) units in the first month and then \(25\) more units each month. With \(u_1=480\) and \(r=25\), we get \(u_n=480+25(n-1)\). Thus \(u_{12}=755\): the shop expects to sell \(755\) units in the twelfth month.
\end{frame}
\begin{frame}{Core Course Notes}
In \(u_3+u_4+\cdots+u_{12}\), there are \(12-3+1=10\) terms, not \(9\) or \(12\).
\end{frame}
\begin{frame}{Core Course Notes}
A sequence \((u_n)\) is \textbf{geometric} when there is a real number \(q\) such that \(u_{n+1}=q u_n\). If the first term is \(u_0\), then \(u_n=u_0q^n\). If the first term is \(u_1\), then \(u_n=u_1q^{n-1}\).
\end{frame}
\begin{frame}{Arithmetic sequences: recognising and modelling}
\begin{itemize}
\item Recognise an additive change.
\item Find the common difference.
\item Define a sequence in context.

\end{itemize}
\end{frame}
\begin{frame}{Arithmetic sequences: recognising and modelling}
\begin{columns}[T,onlytextwidth]
\begin{column}{.57\textwidth}
A club has \(320\) members. It expects to gain \(18\) more each year. Let \(u_n\) be the number of members after \(n\) years. Calculate \(u_0,u_1,u_2\). Which operation takes you from one term to the next? Write a relation between \(u_{n+1}\) and \(u_n\), then predict \(u_5\).
    center center
\end{column}
\begin{column}{.4\textwidth}
\centering
\includegraphics[width=\linewidth,height=4.5cm,keepaspectratio]{figures/stmg-01-3-en.pdf}
\end{column}
\end{columns}
\end{frame}
\begin{frame}{Arithmetic sequences: recognising and modelling}
A sequence is arithmetic when \(u_{n+1}-u_n=r\) is constant. In a real-life situation, a constant absolute increase leads to an arithmetic sequence.
\end{frame}
\begin{frame}{Arithmetic sequences: recognising and modelling}
Take the data from the investigation and write up the solution in this order: definition of the objects, property used, exact or approximate calculation, then a sentence giving the answer. The model answer appears in the teacher edition.
\end{frame}
\begin{frame}{Arithmetic sequences: recognising and modelling}
\textbf{Running a 1-hour session:} 7 min of fluency; 10 min of sharing and written record; 28 min of essential exercises; 5 min of correction; 3 min of exit ticket.

\textbf{Running a 2-hour session:} add 20 min of investigation, 20 min of formalising the rule, a second independent exercise and 10 min of comparative correction.

\textbf{Differentiation:} give struggling students a four-box sheet (data, model, calculation, conclusion). Ask fast finishers to produce a counterexample, a justification or a second method.
\end{frame}
\begin{frame}{Arithmetic sequences: the general term}
\begin{itemize}
\item Move from the recurrence relation to the general term.
\item Calculate a distant term.
\item Interpret the index.

\end{itemize}
\end{frame}
\begin{frame}{Arithmetic sequences: the general term}
A shop has \(1{,}250\) active loyalty cards at index \(0\) and gains \(75\) per month. Without calculating all the terms, find the number of cards at index \(n\), then at index \(18\). Explain the role of \(18\) in the calculation.
\end{frame}
\begin{frame}{Arithmetic sequences: the general term}
If \((u_n)\) is arithmetic with common difference \(r\), then \(u_n=u_0+nr\) or \(u_n=u_1+(n-1)r\), depending on the index of the first term.
\end{frame}
\begin{frame}{Arithmetic sequences: the general term}
Take the data from the investigation and write up the solution in this order: definition of the objects, property used, exact or approximate calculation, then a sentence giving the answer. The model answer appears in the teacher edition.
\end{frame}
\begin{frame}{Arithmetic sequences: the general term}
\textbf{Running a 1-hour session:} 7 min of fluency; 10 min of sharing and written record; 28 min of essential exercises; 5 min of correction; 3 min of exit ticket.

\textbf{Running a 2-hour session:} add 20 min of investigation, 20 min of formalising the rule, a second independent exercise and 10 min of comparative correction.

\textbf{Differentiation:} give struggling students a four-box sheet (data, model, calculation, conclusion). Ask fast finishers to produce a counterexample, a justification or a second method.
\end{frame}
\begin{frame}{Sum of the terms of an arithmetic sequence}
\begin{itemize}
\item Count the terms of a sum.
\item Apply the formula.
\item Check the order of magnitude.

\end{itemize}
\end{frame}
\begin{frame}{Sum of the terms of an arithmetic sequence}
A company pays a monthly bonus: \(45\) euros in the first month, then \(5\) euros more each month. Calculate the total paid over the first twelve months in three ways: explicit sum, pairing first and last terms, formula.
\end{frame}
\begin{frame}{Sum of the terms of an arithmetic sequence}
A sum of arithmetic terms equals \(\text{number of terms}\times\dfrac{\text{first}+\text{last}}2\). From \(u_p\) to \(u_q\), there are \(q-p+1\) terms.
\end{frame}
\begin{frame}{Sum of the terms of an arithmetic sequence}
Take the data from the investigation and write up the solution in this order: definition of the objects, property used, exact or approximate calculation, then a sentence giving the answer. The model answer appears in the teacher edition.
\end{frame}
\begin{frame}{Sum of the terms of an arithmetic sequence}
\textbf{Running a 1-hour session:} 7 min of fluency; 10 min of sharing and written record; 28 min of essential exercises; 5 min of correction; 3 min of exit ticket.

\textbf{Running a 2-hour session:} add 20 min of investigation, 20 min of formalising the rule, a second independent exercise and 10 min of comparative correction.

\textbf{Differentiation:} give struggling students a four-box sheet (data, model, calculation, conclusion). Ask fast finishers to produce a counterexample, a justification or a second method.
\end{frame}
\begin{frame}{Geometric sequences: recognising and modelling}
\begin{itemize}
\item Recognise a multiplicative change.
\item Convert a rate into a multiplier.
\item Compare arithmetic and geometric sequences.

\end{itemize}
\end{frame}
\begin{frame}{Geometric sequences: recognising and modelling}
\begin{columns}[T,onlytextwidth]
\begin{column}{.57\textwidth}
An online shop increases its number of visitors by \(6\%\) per month, starting from \(8{,}000\) visitors. Calculate the first three values. Compare the differences, then the quotients. Identify the type of sequence.
    center center
\end{column}
\begin{column}{.4\textwidth}
\centering
\includegraphics[width=\linewidth,height=4.5cm,keepaspectratio]{figures/stmg-01-4-en.pdf}
\end{column}
\end{columns}
\end{frame}
\begin{frame}{Geometric sequences: recognising and modelling}
A sequence is geometric when \(\dfrac{u_{n+1}}{u_n}=q\) is constant (non-zero terms). A constant relative change leads to a geometric sequence.
\end{frame}
\begin{frame}{Geometric sequences: recognising and modelling}
Take the data from the investigation and write up the solution in this order: definition of the objects, property used, exact or approximate calculation, then a sentence giving the answer. The model answer appears in the teacher edition.
\end{frame}
\begin{frame}{Geometric sequences: recognising and modelling}
\textbf{Running a 1-hour session:} 7 min of fluency; 10 min of sharing and written record; 28 min of essential exercises; 5 min of correction; 3 min of exit ticket.

\textbf{Running a 2-hour session:} add 20 min of investigation, 20 min of formalising the rule, a second independent exercise and 10 min of comparative correction.

\textbf{Differentiation:} give struggling students a four-box sheet (data, model, calculation, conclusion). Ask fast finishers to produce a counterexample, a justification or a second method.
\end{frame}
\begin{frame}{Geometric sequences: general term and sum}
\begin{itemize}
\item Write the general term.
\item Calculate a geometric sum.
\item Interpret a cumulative change.

\end{itemize}
\end{frame}
\begin{frame}{Geometric sequences: general term and sum}
A capital of \(250{,}000\) Ar is invested at \(4\%\) per year, compound interest. Write the capital \(C_n\) after \(n\) years. Calculate \(C_6\), then the sum \(C_0+\cdots+C_6\), and explain what this last sum represents mathematically and why it is not the final capital.
\end{frame}
\begin{frame}{Geometric sequences: general term and sum}
If \((u_n)\) is geometric with common ratio \(q\), \(u_n=u_0q^n\) ou \(u_n=u_1q^{n-1}\). For \(q\neq1\), \(u_0+\cdots+u_n=u_0\dfrac{1-q^{n+1}}{1-q}\).
\end{frame}
\begin{frame}{Geometric sequences: general term and sum}
Take the data from the investigation and write up the solution in this order: definition of the objects, property used, exact or approximate calculation, then a sentence giving the answer. The model answer appears in the teacher edition.
\end{frame}
\begin{frame}{Geometric sequences: general term and sum}
\textbf{Running a 1-hour session:} 7 min of fluency; 10 min of sharing and written record; 28 min of essential exercises; 5 min of correction; 3 min of exit ticket.

\textbf{Running a 2-hour session:} add 20 min of investigation, 20 min of formalising the rule, a second independent exercise and 10 min of comparative correction.

\textbf{Differentiation:} give struggling students a four-box sheet (data, model, calculation, conclusion). Ask fast finishers to produce a counterexample, a justification or a second method.
\end{frame}
\begin{frame}{Review and formative assessment}
\begin{itemize}
\item Choose the right model.
\item Use the formulas.
\item Analyse one's own mistakes.

\end{itemize}
\end{frame}
\begin{frame}{Review and formative assessment}
Sort eight situations into three categories: arithmetic, geometric, neither. For each situation, underline the word or piece of data that justifies the choice, then choose the appropriate formula.
\end{frame}
\begin{frame}{Review and formative assessment}
Before calculating, identify: type of sequence, first index, common difference or ratio, and the term or sum asked for. The conclusion states the unit and the index.
\end{frame}
\begin{frame}{Review and formative assessment}
Take the data from the investigation and write up the solution in this order: definition of the objects, property used, exact or approximate calculation, then a sentence giving the answer. The model answer appears in the teacher edition.
\end{frame}
\begin{frame}{Review and formative assessment}
\textbf{Running a 1-hour session:} 7 min of fluency; 10 min of sharing and written record; 28 min of essential exercises; 5 min of correction; 3 min of exit ticket.

\textbf{Running a 2-hour session:} add 20 min of investigation, 20 min of formalising the rule, a second independent exercise and 10 min of comparative correction.

\textbf{Differentiation:} give struggling students a four-box sheet (data, model, calculation, conclusion). Ask fast finishers to produce a counterexample, a justification or a second method.
\end{frame}
\end{document}
